\section{May 12}


\paragraph{Reference:} [Bloch:``algebraic cycles and higher K-theory''], [Bloch:``The moving lemma for higher chow groups''] [Mazza-Voevodsky-Weibel:``lecture notes on motivic cohomology'']


\subsection{Higher Chow groups}

  For a closed subscheme $Z$ of a scheme $X$, we have the localization sequence
  \[ \CH_i(Z) \to \CH_i(X) \to \CH_i(X-Z) \to 0, \]
  which is only exact on the right. 
  Similarly, for open subscheme $U_1,U_2$ of $X$, we have the Mayer-Vietoris sequence
  \[ \CH_i(X) \to \CH_i(U_1) \oplus \CH_i(U_2) \to \CH_i(U_1\cap  U_2) \to 0, \]
  which is also only exact on the right.
  For example, let $X = \bbP^1$, $U_1 = \bbP^1 - \{0\}$, $U_2 = \bbP^1 - \{\infty\}$, then $\CH_0(X) = \bbZ$, $\CH_0(U_1) = \CH_0(U_2) = 0$, $\CH_0(U_1\cap U_2) = \bbZ$.
  
  Conceptually, we want a theory of higher Chow groups $\CH_i(X,n)$, so that 
  \[ K_n(X)_\bbQ \cong \bigoplus_i \CH_i(X,n)_\bbQ, \]
  where $K_n(X)_\bbQ$ is the $n$-th algebraic K-group of $X$, defined as the higher homotopy groups of a certain space associated to $X$, for which the localization sequence is exact on the left by topological reasons.
  We will construct $\CH_i(X,n)$ using a complex of abelian groups, modeled upon singular complex.

  Define the algebraic $n$-simplex $\Delta^n := \Spec k[t_0,\dots,t_n]/(\sum t_i - 1) \cong \bbA^n$.
  For an increasing map $\rho: [m] \to [n]$, define $\tilde{\rho}: \Delta^m \to \Delta^n$ by $\tilde{\rho}^*(t_i) = \sum_{j\in \rho^{-1}(i)} t_j$.
  Here $[n]$ is the set $\{0,1,\dots,n\}$.
  If $\rho$ is injective, we say $\tilde{\rho}(\Delta^m)$ is a face of $\Delta^n$.
  
  Let $X$ be an equi-dimensional scheme (separated and of finite type over $\kk$).
  Define $Z^i(X,n) \subset Z^i(X\times \Delta^n)$ to be the subgroup generated by all codimension $i$ subvarieties of $X\times \Delta^n$ which intersect all faces of $X\times \Delta^n$ properly.
  For instance, if $n=2$ and $X$ is a point, then $Z^1(X,2)$ is generated by all lines in $\Delta^2 \subset \bbA^3$ which do not pass through any vertex of $\Delta^2$.\Yang{}

  For each $l \in [n]$, define $\partial_l: Z^i(X,n) \to Z^i(X,n-1)$ by the refined Gysin pull-back along the face $\tilde{\rho}_l: \Delta^{n-1} \to \Delta^n$ corresponding to the injective map $\rho_l: [n-1] \to [n]$ which misses $l$.
  It is well-defined because of the proper intersection condition.
  We get a complex of abelian groups:
  \[ Z^i(X,\bullet): \cdots \to Z^i(X,2) \xrightarrow{\sum_{l=0}^2 (-1)^l \partial_l} Z^i(X,1) \xrightarrow{ \sum_{l=0}^1 (-1)^l \partial_l} Z^i(X,0) \to 0. \]

  \begin{definition}
    The $n$-th higher Chow group of $X$ is defined as the $n$-th homology of the complex $Z^i(X,\bullet)$.
  \end{definition}

  The higher Chow groups vanish for $n<0$ automatically.

  \begin{example}
    When $n=0$, the proper intersection condition is vacuous, so $Z^i(X,0) = Z^i(X)$, and 
    \[ \CH^i(X,0) = \coker(Z^i(X,1) \to Z^i(X,0)) = Z^i(X)/\Image(Z^i(X,1) \to Z^i(X,0)) = \CH^i(X) \]
    since $\Image(Z^i(X,1) \to Z^i(X,0))$ is generated by all cycles of the form $[Z_0] - [Z_1]$ where $Z$ is a codimension $i$ subvariety of $X\times \Delta^1$ which intersects $X\times \{0\}$ and $X\times \{1\}$ properly, and $Z_0,Z_1$ are the intersections of $Z$ with $X\times \{0\}$ and $X\times \{1\}$ respectively.
  \end{example}


\subsection[Properties of higher Chow groups]{Properties of $\CH^i(X,n)$}

  In the following, all schemes are assumed to be equi-dimensional and quasi projective.


  \begin{proposition}
    \begin{enumerate}
      \item: Functoriality: 
        \begin{enumerate}
          \item let $f: X \to Y$ be a proper morphism with $e = \dim X - \dim Y$, then we have a push-forward map $f_*: \CH^i(X,n) \to \CH^{i-e}(Y,n)$;
          \item let $f: X \to Y$ be a flat morphism, then we have a pull-back map $f^*: \CH^i(Y,n) \to \CH^i(X,n)$;
          \item let $f: X \to Y$ be a morphism, where $Y$ is smooth, then we have a pull-back map $f^*: \CH^i(Y,n) \to \CH^i(X,n)$.
        \end{enumerate}
      \item Homotopy invariance: let $\pi: E \to X$ be a vector bundle, then we have $\pi^*: \CH^i(X,n) \to \CH^i(E,n)$ is an isomorphism.
      \item Localization: let $Y \subset X$ be a closed subscheme of pure codimension $d$, then we have a long exact sequence
        \[ \cdots \to \CH^{i-d}(Y,n) \to \CH^i(X,n) \to \CH^i(X-Y,n) \to \CH^{i-d}(Y,n-1) \to \cdots. \]
      \item Products: we have a product map
        \[ \CH^i(X,n) \otimes \CH^j(Y,m) \to \CH^{i+j}(X\times Y, n+m). \]
        When $X$ is smooth, pullback along the diagonal $\Delta: X \to X\times X$ gives a product map
        \[ \CH^i(X,n) \otimes \CH^j(X,m) \to \CH^{i+j}(X, n+m). \]
      \item Chern classes: let $E \to X$ be a vector bundle of rank $r$, then for all $j=0,1,\dots,r$, we have a Chern class $c_j(E): \CH^i(X,n) \to \CH^{i+j}(X,n)$, letting $\xi = c_1(\calO_{\bbP(E)}(1)) \in \CH^1(\bbP(E),0)$, we 
      \item Relation with algebraic K-theory: when $X$ is smooth, we have an isomorphism
        \[ K_n^Q(X)_\bbQ \cong \bigoplus_i \CH^i(X,n)_\bbQ. \]
        Here $K_n^Q(X)$ is the $n$-th Quillen K-group of $X$ on the category of vector bundles on $X$.
    \end{enumerate}
  \end{proposition}
  \begin{remark}
    The proof of (c) in [B1] contains an error, which was corrected in [B2].
    The proof of (a.iii) in [B1] contains an error, which was fixed by Levine (correct proof in [MVW]).
  \end{remark}
  \begin{proof}
    \Yang{}
  \end{proof}

  \begin{lemma}
    Let $X$ be a scheme, $G$ be an algebraic group acting on $X$.
    Let $A,B \subset X$ be closed subsets.
    Assume $G\times A \to X$ satisfies
    \begin{itemize}
      \item its fibers have the same dimension;
      \item it is dominant onto $X$.
    \end{itemize}
    Then there exists $U \subset G$ non-empty open, s.t. for all $g \in U$, $gA$ intersects $B$ properly.
  \end{lemma}
  \begin{proof}
    \Yang{}
  \end{proof}

  \begin{lemma}
    Under the assumptions of the previous lemma, we more assume 
    \begin{itemize}
      \item there exists $\KK/\kk$ extension,
      \item there exists $\psi: X_\KK \to G_\KK$ a morphism
      \item there exists $V \subset X$ non-empty open
    \end{itemize}
    such that for all $x \in V_\KK$, $\trdeg_\kk \kk(\pi(x), \pi'(\psi(x))) \geq \dim G$, where $\pi: X_\KK \to X$ and $\pi': G_\KK \to G$ are the projections.
    Define $\phi: X_\KK \to X_\KK, x \mapsto \psi(x)x$.
    Assume $\phi$ is an isomorphism.
    Then $\phi(A_\kk \cap V)$ and $B_\KK$ intersect properly.
  \end{lemma}
  \begin{proof}
    \Yang{}
  \end{proof}